\documentclass[border=10pt]{standalone}
\usepackage[no-math]{fontspec}
\usepackage[brazil]{babel}
\usepackage{amsmath,amssymb}
\IfFileExists{C:/Windows/Fonts/Lato_Regular.ttf}{%
  \setmainfont{Lato}[
    Path = C:/Windows/Fonts/,
    UprightFont = Lato_Regular.ttf,
    BoldFont = Lato_Bold.ttf,
    ItalicFont = Lato_Italic.ttf,
    BoldItalicFont = Lato_BoldItalic.ttf
  ]
  \setsansfont{Lato}[
    Path = C:/Windows/Fonts/,
    UprightFont = Lato_Regular.ttf,
    BoldFont = Lato_Bold.ttf,
    ItalicFont = Lato_Italic.ttf,
    BoldItalicFont = Lato_BoldItalic.ttf
  ]
}{%
  \RequirePackage[default]{lato}
}
\renewcommand{\familydefault}{\sfdefault}

\usepackage{xcolor}
\usepackage{tikz}
\usepackage{pgfplots}
\pgfplotsset{compat=1.18}
\usetikzlibrary{arrows.meta,positioning,shapes.geometric,fit,backgrounds,calc,decorations.pathreplacing,patterns}

% Paleta institucional UEFS / Extensão Rural
\definecolor{erNavy}{HTML}{1F4268}
\definecolor{erNavyDark}{HTML}{152E4A}
\definecolor{erGold}{HTML}{B38A3D}
\definecolor{erGoldDark}{HTML}{8C6927}
\definecolor{erTeal}{HTML}{0E7C7B}
\definecolor{erForest}{HTML}{2D6A4F}
\definecolor{erAlert}{HTML}{9E2A2B}
\definecolor{erInk}{HTML}{26313D}
\definecolor{erMuted}{HTML}{5B6573}

\definecolor{bgNavy}{HTML}{EDF2F7}
\definecolor{bgGold}{HTML}{F8F3E9}
\definecolor{bgTeal}{HTML}{EBF6F5}
\definecolor{bgForest}{HTML}{EEF5F0}
\definecolor{bgAlert}{HTML}{FDF2F2}

\tikzset{
  >=Stealth,
  erBox/.style={
    draw=#1, fill=white, rounded corners=5pt, line width=1.15pt,
    inner sep=7pt, align=center, font=\footnotesize\sffamily, text=erInk
  },
  erTint/.style={
    draw=#1, fill=#1!8!white, rounded corners=5pt, line width=1.15pt,
    inner sep=7pt, align=center, font=\footnotesize\sffamily, text=erInk
  },
  erBadge/.style={
    fill=#1, text=white, rounded corners=9pt, inner xsep=8pt, inner ysep=3pt,
    font=\scriptsize\bfseries\sffamily
  },
  erTitle/.style={font=\normalsize\bfseries\sffamily, text=erNavy, align=center},
  erSub/.style={font=\footnotesize\sffamily, text=erMuted, align=center},
  erArrow/.style={->, >=Stealth, line width=1.35pt},
  erDArrow/.style={<->, >=Stealth, line width=1.35pt}
}


\begin{document}
\begin{tikzpicture}[font=\sffamily]

% Título e subtítulo bem acima do degrau 4
\node[erTitle, font=\large\bfseries\sffamily] (title) at (8.05, 7.35)
  {Os 4 Níveis da Transição Agroecológica (Gliessman, 2001; 2015)};
\node[erSub] (sub) at (8.05, 6.85)
  {Escalonamento progressivo do manejo técnico à transformação socioecológica do território};

% Cards em degraus ascendentes (esquerda -> direita) com espaçamento para setas
\node[erTint=erMuted, text width=3.15cm, minimum height=2.25cm, anchor=south west] (n1) at (0.0, 0.95) {
  \textbf{\textcolor{erInk}{Nível 1}}\\[2pt]
  \textbf{\textcolor{erNavy}{Eficiência de Insumos}}\\[5pt]
  \scriptsize\textit{Reduzir desperdício e doses de agroquímicos}
};

\node[erTint=erGold, text width=3.15cm, minimum height=2.25cm, anchor=south west] (n2) at (4.20, 1.85) {
  \textbf{\textcolor{erGoldDark}{Nível 2}}\\[2pt]
  \textbf{\textcolor{erNavy}{Substituição de Insumos}}\\[5pt]
  \scriptsize\textit{Bioinsumos, adubação verde, controle biológico}
};

\node[erTint=erTeal, text width=3.15cm, minimum height=2.25cm, anchor=south west] (n3) at (8.40, 2.75) {
  \textbf{\textcolor{erTeal}{Nível 3}}\\[2pt]
  \textbf{\textcolor{erNavy}{Redesenho do Agroecossistema}}\\[5pt]
  \scriptsize\textit{Policultivos, SAFs,\\ciclagem de nutrientes}
};

\node[erTint=erForest, text width=3.15cm, minimum height=2.25cm, anchor=south west] (n4) at (12.60, 3.65) {
  \textbf{\textcolor{erForest}{Nível 4}}\\[2pt]
  \textbf{\textcolor{erNavy}{Conexão Direta \& Territorialidade}}\\[5pt]
  \scriptsize\textit{Feiras, SPGs, cultura alimentar sustentável}
};

% Badges de escopo acima de cada degrau (rounded corners ajustado)
\node[erBadge=erMuted, rounded corners=5pt, anchor=south] at ([yshift=5pt]n1.north) {Racionalização};
\node[erBadge=erGold, rounded corners=5pt, anchor=south] at ([yshift=5pt]n2.north) {Substituição};
\node[erBadge=erTeal, rounded corners=5pt, anchor=south] at ([yshift=5pt]n3.north) {Redesenho};
\node[erBadge=erForest, rounded corners=5pt, anchor=south] at ([yshift=5pt]n4.north) {Sistema Agroalimentar};

% Setas de progressão entre degraus
\draw[erArrow, color=erGoldDark] ([yshift=14pt]n1.east) -- ([yshift=-14pt]n2.west);
\draw[erArrow, color=erTeal] ([yshift=14pt]n2.east) -- ([yshift=-14pt]n3.west);
\draw[erArrow, color=erForest] ([yshift=14pt]n3.east) -- ([yshift=-14pt]n4.west);

% Linha base e seta direcional inferior
\draw[erArrow, color=erNavy, line width=2pt] (0.0, 0.30) -- (16.10, 0.30);
\node[fill=white, inner xsep=10pt, font=\footnotesize\bfseries\sffamily, text=erNavy]
  at (8.05, 0.30) {Da Escala da Parcela à Transformação do Sistema Agroalimentar $\rightarrow$};

\end{tikzpicture}
\end{document}
